% !TEX root = ../main.tex
\chapter{Data Loading Pipeline}
\label{ch:data_pipeline}

The fastest possible model is useless if it sits idle waiting
for data. This chapter treats \Lucid's data-loading
infrastructure: the \code{Dataset} and \code{DataLoader}
classes, the worker-process pool that runs CPU-side
preprocessing in parallel, the prefetch queue that keeps the
GPU fed, the standard transform library
(image, text, audio), the data-augmentation primitives, and the
performance pitfalls that hurt the worst (and how to spot
them). The data pipeline is the \emph{only} compute path in
\Lucid\ that interacts with the outside world (file I/O, decoded
images, etc.), so it is also the chapter where the
``bridge-boundary'' rule \cite{numpy_van_der_walt_2011}
matters most.

\section{The Dataset / DataLoader Pattern}
\label{sec:data_pipeline:dataset_dataloader}

The standard pattern, inherited from torchvision conventions:

\begin{lstlisting}[style=lucid,language=LucidPython]
class Dataset:
    def __len__(self) -> int: ...
    def __getitem__(self, idx: int) -> tuple[Tensor, Tensor]: ...

dataset = MyDataset(...)
loader = DataLoader(
    dataset,
    batch_size=32,
    shuffle=True,
    num_workers=4,
    pin_memory=True,
)
for x, y in loader:
    out = model(x)
    loss = loss_fn(out, y)
    loss.backward()
\end{lstlisting}

The \code{Dataset} is a per-item interface; the
\code{DataLoader} batches, shuffles, parallelises, and prefetches.

\subsection{Why the split}
\label{ssec:data_pipeline:why_split}

The dataset's per-item interface is decoupled from the loader's
batching and parallelism concerns. This lets the same dataset
work with multiple loaders (e.g., one for training with
augmentation and shuffling, one for evaluation without), and it
lets the loader be replaced (e.g., with a distributed version)
without touching the dataset.

\section{Worker Process Pool}
\label{sec:data_pipeline:worker_pool}

CPU-side preprocessing (image decode, augmentation, tokenisation)
is parallelised across worker processes. The \code{num\_workers}
parameter controls the pool size; common values are
$N_{\text{CPU}} / 2$ for a balance between data prep and model
compute.

\subsection{Why processes, not threads}
\label{ssec:data_pipeline:processes_vs_threads}

Python's GIL prevents true thread parallelism for CPU-bound
work. The standard pattern is multiprocessing: each worker is
a separate Python process with its own GIL, communicating with
the main process via shared memory.

\subsection{Inter-process communication}
\label{ssec:data_pipeline:ipc}

Workers send produced batches back to the main process via:
\begin{itemize}
  \item POSIX shared memory for tensor data (zero copy on
    Apple silicon, where all memory is unified).
  \item A small pipe for metadata (item indices, exception
    info).
\end{itemize}

The combination achieves $<$\,1\,ms IPC overhead per batch on
typical workloads, well below the model-compute time.

\subsection{Worker initialisation}
\label{ssec:data_pipeline:worker_init}

Each worker process needs to:
\begin{itemize}
  \item Import \Lucid\ (about 400\,ms; amortised over many
    batches).
  \item Set a per-worker random seed (so each worker generates
    different augmentations).
  \item Open per-worker file handles (e.g., LMDB, HDF5).
\end{itemize}

The \code{worker\_init\_fn} parameter lets the user customise:

\begin{lstlisting}[style=lucid,language=LucidPython]
def worker_init_fn(worker_id):
    lucid.seed(42 + worker_id)
    open_lmdb_for_worker(worker_id)

loader = DataLoader(..., worker_init_fn=worker_init_fn)
\end{lstlisting}

\section{Prefetch Queue}
\label{sec:data_pipeline:prefetch_queue}

The data loader's prefetch queue maintains a buffer of ready
batches ahead of the consumer. The default queue size is
\code{2 * num\_workers}: enough to absorb short-term
variation in worker latency without burning excessive memory.

\begin{lstlisting}[style=lucid,language=LucidPython]
loader = DataLoader(
    dataset,
    batch_size=32,
    num_workers=4,
    prefetch_factor=2,  # 2 * num_workers = 8 batches buffered
)
\end{lstlisting}

For workloads where the GPU step is fast (small models), a
small prefetch is enough. For workloads with expensive per-item
work (decoding video frames, complex augmentations), a larger
prefetch may be needed.

\section{Pin Memory}
\label{sec:data_pipeline:pin_memory}

In CUDA-based frameworks, ``pinned'' memory is a CPU page that
is page-locked, allowing direct DMA to the GPU without an
intermediate copy. On Apple silicon, unified memory means there
is no CPU--GPU copy at all; \code{pin\_memory=True} is a no-op
(but accepted for compatibility).

\section{Standard Transforms}
\label{sec:data_pipeline:standard_transforms}

\Lucid\ provides a \code{lucid.utils.data.transforms} module
covering the common image transforms:

\begin{itemize}
  \item \code{Resize(size)}: resize via Pillow or
    \code{cv2}.
  \item \code{RandomCrop(size)}: random crop with optional
    padding.
  \item \code{RandomHorizontalFlip(p=0.5)}.
  \item \code{ColorJitter(brightness, contrast, saturation,
    hue)}.
  \item \code{Normalize(mean, std)}: per-channel
    standardisation.
  \item \code{ToTensor()}: PIL Image $\to$ \Lucid\ Tensor.
\end{itemize}

Composition via \code{Compose}:

\begin{lstlisting}[style=lucid,language=LucidPython]
from lucid.utils.data import transforms as T

train_transform = T.Compose([
    T.RandomResizedCrop(224),
    T.RandomHorizontalFlip(),
    T.ColorJitter(0.4, 0.4, 0.4),
    T.ToTensor(),
    T.Normalize(mean=[0.485, 0.456, 0.406], std=[0.229, 0.224, 0.225]),
])
\end{lstlisting}

The transforms work on PIL or NumPy inputs; the
\code{ToTensor} step is the boundary to \Lucid\ tensor.

\section{Standard Augmentation Recipes}
\label{sec:data_pipeline:augmentation_recipes}

For convenience, \Lucid\ also ships pre-configured recipes:

\begin{lstlisting}[style=lucid,language=LucidPython]
from lucid.utils.data.augment import imagenet_train_transform

transform = imagenet_train_transform()
# Encapsulates: RandomResizedCrop(224) + RandomHorizontalFlip +
#   AugMix + ColorJitter(0.4, 0.4, 0.4) + ToTensor + ImageNet normalise
\end{lstlisting}

The recipes are versioned: \code{imagenet\_train\_transform\_v2}
adds the recipe used by recent state-of-the-art models. Old
versions remain available for reproducibility of older results.

\section{Collation}
\label{sec:data_pipeline:collation}

The default \code{collate\_fn} batches a list of samples into
a single batch tensor. For uniform shapes (image classification),
this is straightforward stacking. For variable-shape samples
(detection, language modelling with variable-length sequences),
a custom collate\_fn pads:

\begin{lstlisting}[style=lucid,language=LucidPython]
def detection_collate(samples):
    images = lucid.stack([s["image"] for s in samples])
    # Detection targets are variable-length per image:
    targets = [{
        "boxes": s["boxes"],
        "labels": s["labels"],
    } for s in samples]
    return images, targets

loader = DataLoader(..., collate_fn=detection_collate)
\end{lstlisting}

For language modelling, the collate pads sequences to the
batch's maximum length and computes the attention mask:

\begin{lstlisting}[style=lucid,language=LucidPython]
def lm_collate(samples):
    max_len = max(s.size(0) for s in samples)
    padded = lucid.zeros(len(samples), max_len, dtype=lucid.int64)
    mask = lucid.zeros(len(samples), max_len, dtype=lucid.bool)
    for i, s in enumerate(samples):
        padded[i, :s.size(0)] = s
        mask[i, :s.size(0)] = True
    return padded, mask
\end{lstlisting}

\section{Performance Diagnostics}
\label{sec:data_pipeline:diagnostics}

The most common training-throughput issue is the data pipeline
under-feeding the model. Symptoms:
\begin{itemize}
  \item GPU utilisation $<$\,80\% during steady-state training.
  \item \code{lucid.profiler} attributes time to ``waiting for
    next batch''.
  \item \code{nvtop}-equivalent monitor shows GPU idle gaps.
\end{itemize}

\Lucid's diagnostic helper:

\begin{lstlisting}[style=lucid,language=LucidPython]
from lucid.utils.data import diagnose

diag = diagnose(loader, num_batches=50)
print(diag.summary())
# DataLoader diagnostic:
#   Median batch latency: 12.5 ms
#   p95 batch latency: 38.2 ms (concerning)
#   Workers seem under-loaded:
#     - num_workers=4, but only 60% utilization observed
#     - suggest reducing to 3 workers OR per-item work too uneven
\end{lstlisting}

The diagnostic measures the queue's fill level over time; a
chronically-empty queue indicates worker under-provisioning.

\section{Common Pitfalls}
\label{sec:data_pipeline:pitfalls}

\subsection{Worker startup cost}
\label{ssec:data_pipeline:worker_startup_cost}

Each worker process imports \Lucid\ (\(\sim\)\,400\,ms). For
short training runs (a few minutes), this is a substantial
fraction of total time. For very short runs, prefer
\code{num\_workers=0} (in-process loading).

\subsection{Per-item python overhead}
\label{ssec:data_pipeline:py_overhead}

Pure-Python per-item code (e.g., complex augmentation logic) is
slow. If profiling shows the workers as the bottleneck,
candidate fixes:
\begin{itemize}
  \item Cache the decoded data (decode once at dataset init).
  \item Move expensive logic to a compiled path (Cython,
    Pillow's C decoders, etc.).
  \item Move augmentation to GPU via
    \code{lucid.gpu\_augment} (where the augmentation is
    representable as tensor ops).
\end{itemize}

\subsection{File-system bottleneck}
\label{ssec:data_pipeline:fs_bottleneck}

Reading millions of small files is slow on most filesystems.
For large datasets, prefer pre-packed formats:
\begin{itemize}
  \item LMDB or LevelDB for key-value access.
  \item Tar shards (WebDataset-style) for streaming.
  \item HDF5 for typed array storage.
\end{itemize}

\Lucid\ provides reader classes for each.

\subsection{Spurious synchronisation}
\label{ssec:data_pipeline:spurious_sync}

The data loader receives label tensors; the natural pattern
\code{label.item()} (to convert to a Python int) forces a
device-to-CPU sync. The
\chref{ch:stabilization}~Section~\ref{ssec:stabilization:phase_4_fixes}
fix changed the loader to keep labels as tensors throughout.

\section{Multi-Process Determinism}
\label{sec:data_pipeline:determinism}

For reproducible training, the data order and augmentations
must be deterministic across runs. \Lucid's
\code{DataLoader} supports this via:

\begin{lstlisting}[style=lucid,language=LucidPython]
loader = DataLoader(
    dataset,
    batch_size=32,
    num_workers=4,
    generator=lucid.random.manual_seed(42),
    worker_init_fn=lambda wid: lucid.seed(42 + wid),
)
\end{lstlisting}

The \code{generator} controls the shuffling order; the
\code{worker\_init\_fn} controls the per-worker augmentation
seed. Together, two runs with the same seed produce the same
data sequence.

\section{Distributed Data Loading}
\label{sec:data_pipeline:distributed}

For multi-node training (not yet a 3.5 feature, but planned),
the data sampler distributes items across nodes:

\begin{lstlisting}[style=lucid,language=LucidPython]
from lucid.utils.data import DistributedSampler

sampler = DistributedSampler(
    dataset, num_replicas=world_size, rank=local_rank,
)
loader = DataLoader(dataset, batch_size=32, sampler=sampler)
\end{lstlisting}

Each rank sees a disjoint subset of the dataset. The sampler
also handles per-epoch shuffling across all ranks.

\section{Tokenisation in the Pipeline}
\label{sec:data_pipeline:tokenisation_pipeline}

For language workloads, the tokenisation step is a non-trivial
fraction of CPU work. Two strategies:

\begin{enumerate}
  \item \textbf{Pre-tokenise once} (recommended for fixed
    dataset): cache tokenised IDs to disk; the dataset reads
    pre-tokenised IDs at item time. Faster but uses more disk.
  \item \textbf{Tokenise on the fly}: the dataset returns text;
    the collate\_fn tokenises. Slower per item but no pre-cache
    needed.
\end{enumerate}

\Lucid's tokenisers (\code{lucid.models.tokenizers}) are
fast Python implementations (\(\sim\)\,2 \textmu s/token);
for very large datasets, the pre-tokenisation path is still
better.

\section{Image Decoding}
\label{sec:data_pipeline:image_decoding}

Image decoding (JPEG, PNG, WebP) is dominated by the CPU
decoder library. \Lucid\ uses Pillow's libjpeg-turbo backend
by default. For batched decode, the
\code{turbojpeg} acceleration option can be enabled:

\begin{lstlisting}[style=lucid,language=LucidPython]
from lucid.utils.data import ImageFolder

dataset = ImageFolder(
    "imagenet/train/",
    transform=train_transform,
    decoder="turbojpeg",  # ~2x faster than default Pillow
)
\end{lstlisting}

For HEIC/AVIF images (common on Apple platforms), the decoder
falls back to the system's image-I/O framework.

\section{Performance Benchmark}
\label{sec:data_pipeline:perf_bench}

ImageFolder + ResNet-50 training, M3~Max, batch 32:

\begin{table}[t]
\centering
\small
\begin{tabular}{lr}
\toprule
num\_workers & End-to-end iter (ms) \\
\midrule
0  (in-process) & 215 \\
1  & 165 \\
2  & 105 \\
4  & 78  \\
8  & 72  \\
16 & 76  \\
\bottomrule
\end{tabular}
\caption[Worker count sweep.]{Iteration time as
\code{num\_workers} varies. Below 4 workers, the loader is the
bottleneck. Above 8, diminishing returns (the data prep is no
longer the limit). Optimal is 4--8 on M3~Max for typical
ImageNet workloads.}
\label{tab:data_pipeline:worker_sweep}
\end{table}

The optimal \code{num\_workers} depends on the per-item work
and the model's compute time. \Lucid's diagnostic helper
suggests a value based on observed queue behaviour.

\section{Summary and Forward References}
\label{sec:data_pipeline:summary}

The data pipeline: Dataset + DataLoader pattern, multi-process
workers for CPU prep, prefetch queue to keep the GPU fed,
standard image-augmentation transforms, custom collate\_fns
for variable-shape inputs, diagnostics for spotting
bottlenecks. The pipeline is the one place \Lucid\ interacts
with the outside-world file system (per Rule H4), so it
includes the necessary bridge code to NumPy / PIL / etc.

The next chapter, \chref{ch:error_handling}, treats error
messaging. \chref{ch:reproducibility} treats reproducibility.

\begin{lucidnote}
For the user: start with \code{num\_workers=4} as a default;
profile if iteration time is dominated by waiting for data.
Use the diagnostic helper (\code{lucid.utils.data.diagnose})
to confirm.
\end{lucidnote}
